% ============================================================
%  theory_model.tex -- companion theory model
%  Spec: docs/superpowers/specs/2026-09-17-toy-model-design.md
%
%  Build from docs/model/:
%    pdflatex theory_model
%    bibtex   theory_model
%    pdflatex theory_model
%    pdflatex theory_model
% ============================================================
\documentclass[11pt]{article}
\usepackage[margin=1.1in]{geometry}
\usepackage{amsmath,amssymb,amsthm}
\usepackage{booktabs}
\usepackage[authoryear,round]{natbib}
\usepackage[hidelinks]{hyperref}

\newtheorem{proposition}{Proposition}
\theoremstyle{remark}
\newtheorem*{remark}{Remark}

\title{What a Filing Reveals:\\ A Companion Model for the Electoral Costs of
Judicialization}
\author{}
\date{\today}

\begin{document}
\maketitle

\begin{abstract}
\noindent
Adversarial electoral litigation widens the top-two margin in Brazilian
mayoral races. This note asks what mechanism can produce that, and finds that
the natural candidate cannot. In the transfer class of campaign-attack models,
an attack moves support from its target to its sender; under the attack
direction those models derive and document, the map is a contraction on the
top-two margin, so it cannot generate the result. The alternative developed
here turns on what a filing reveals. A filing is a noisy signal about a
candidate's type and a precise signal about a candidate's rank, because the
filer is optimizing and targets the leading opponent. Litigation therefore
informs voters about position more than about quality. Coordination on the
revealed leader widens the margin; the collapse of pivot probability sends
some of the runner-up's support off the valid ballot, which under compulsory
voting means a blank or null vote. Whether that support joins the leader or
leaves depends on how settled the leader's reputation is, which is what
separates open from contested seats.
\end{abstract}

\section{Environment}
\label{sec:env}

One mayoral race. Candidates $j = 1,\dots,J$ are indexed by pre-shock support
$s^\circ_1 \ge s^\circ_2 \ge \dots \ge s^\circ_J$, and there is a continuum of
voters of mass one. Voter $i$'s payoff from candidate $j$ is
\[
  u_{ij} \;=\; \hat\theta_j \;-\; \tfrac{\rho}{2}\,\hat V_i(\theta_j)
          \;+\; \varepsilon_{ij},
\]
where $\theta_j$ is a common-value valence nobody observes and
$\varepsilon_{ij}$ is a fixed idiosyncratic taste the voter knows. Beliefs
about valence start at $\theta_j \sim N(m_j, 1/h_j)$: $m_j$ is the candidate's
reputation and $h_j$ measures how settled the record is. An incumbent who has
governed for four years has a high $h$; a first-time candidate has a low one.
Let $\bar b$ denote the payoff from a blank or null vote, which is what
refusing to endorse anyone is worth once the ballot carries no instrumental
value.

Treatment is $\lambda$, the intensity of adversarial filings in the race, and
it is the object the shift-share instrument shifts. Two institutional facts
fix the timing and the margin of adjustment. Candidate registration closes on
15 August while adversarial filings run through October, so the candidate set
is predetermined with respect to $\lambda$. And voting is compulsory
(CF/88 art.~14 \S1), so a voter who wants out of the contest adjusts the mark
rather than the attendance: the exit option is a blank or null vote.

\section{The transfer channel cannot produce the result}
\label{sec:foil}

\subsection{The map}

\citet{nakaguma2025} model a campaign attack as a transfer. An attack from $k$
on $j$ moves a fraction $\phi$ of $j$'s support away from $j$; in their
Appendix G.2 demobilization extension a fraction $\alpha$ of that mass leaves
for the outside option instead of reaching $k$, and the attacker loses a
fraction $\gamma$ of its own support to backlash.

Their equilibrium has every candidate target the highest-ranked opponent, so
in a three-candidate race the profile is $1 \to 2$, $2 \to 1$ and $3 \to 1$:
mutual between the top two, with the third attacking the leader. Their data
agrees. Measuring attacks as electoral lawsuits, the observed frequencies in
three-candidate races are $2 \to 1$ at 3.0pp, $1 \to 2$ at 1.9pp and
$3 \to 1$ at 0.8pp. The second-place candidate is the most aggressive and the
front-runner is the most attacked.

\subsection{The contraction}

\begin{proposition}[Transfer cannot consolidate]
\label{prop:foil}
Under mutual $1 \leftrightarrow 2$ attacks, post-attack support is
\[
  s_1 = s^\circ_1(1 - \phi - \gamma) + \phi(1-\alpha)\,s^\circ_2,
  \qquad
  s_2 = s^\circ_2(1 - \phi - \gamma) + \phi(1-\alpha)\,s^\circ_1,
\]
so the top-two margin satisfies
\[
  s_1 - s_2 \;=\; (s^\circ_1 - s^\circ_2)\,\kappa,
  \qquad
  \kappa \;\equiv\; 1 - \gamma - \phi(2-\alpha).
\]
For every $\phi \in (0,1)$, $\alpha \in [0,1]$ and $\gamma \in [0,1)$ we have
$\kappa < 1$. Adding $3 \to 1$ subtracts a further $\phi\,s^\circ_1$ from the
margin. Hence the transfer map cannot widen the top-two margin while
preserving the identity of the leader.
\end{proposition}

\begin{proof}
Subtracting,
\[
  s_1 - s_2
  = (s^\circ_1 - s^\circ_2)(1 - \phi - \gamma)
    - \phi(1-\alpha)(s^\circ_1 - s^\circ_2)
  = (s^\circ_1 - s^\circ_2)\bigl[1 - \gamma - \phi(2-\alpha)\bigr].
\]
Since $\alpha \le 1$ we have $2 - \alpha \ge 1$, so $\phi(2-\alpha) \ge \phi >
0$, and with $\gamma \ge 0$ this gives $\kappa < 1$. If $\kappa \in (0,1)$ the
margin strictly narrows. If $\kappa \le 0$ the sign of $s_1 - s_2$ is
non-positive, so candidate 2 has caught or passed candidate 1 and the leader
is no longer the pre-attack leader. When $3 \to 1$ is added, candidate 1 loses
a further $\phi\,s^\circ_1$ and candidate 2 loses nothing, so
$s_1 - s_2 = (s^\circ_1 - s^\circ_2)\kappa - \phi\,s^\circ_1$, which is
smaller still.
\end{proof}

\begin{remark}
Widening the margin under the transfer map requires a one-sided downward
attack, $1 \to 2$ with no reciprocation. That is excluded by the equilibrium,
in which candidate 2's best target is candidate 1, and it is not what the data
shows. The per-pair attack frequencies are also 1--3pp, an order of magnitude
too small to move municipality-level vote shares by three points.
\end{remark}

\begin{remark}[Object of comparison]
\citet{nakaguma2025} work with a win probability through a logit contest
function, whereas the estimates here are vote-share margins. The two are
monotone in $s_1 - s_2$, so Proposition~\ref{prop:foil} transfers, but the
mapping is an assumption and not an identity.
\end{remark}

\begin{remark}[What Proposition~\ref{prop:foil} does not do]
\label{rem:orderstat}
The top-two margin is an order statistic, so any mean-preserving increase in
the dispersion of vote shares widens it mechanically.
Proposition~\ref{prop:foil} rules out the transfer channel under the
documented attack direction; it does not rule out a generic dispersion shock.
What rules that out is the ballot response and the contrast between seat
types in Section~\ref{sec:seat}: a dispersion shock predicts neither, and
predicts no directional flow from the runner-up to the leader.
\end{remark}

\section{What a filing emits}
\label{sec:signals}

A filing carries two signals, and their precisions differ by an order of
magnitude.

\paragraph{The type signal.}
A filing is an allegation about $\theta_j$. It is informative with probability
$\pi$ and pure noise with probability $1 - \pi$, and the voter cannot tell
which. Writing $\tau_I$ for the precision of an informative filing, the
effective precision is $\bar\tau = \pi\tau_I$, and the posterior mean is
\[
  \hat\theta_j \;=\; \frac{h_j m_j + \bar\tau s_j}{h_j + \bar\tau}.
\]
Genuine information about a candidate and a weapon that happens to persuade
enter through the same $\bar\tau$. They are therefore not separable for the
voter, and not separable for the econometrician either. The
non-separability that the empirical design has to live with is a property of
the environment, not a shortcoming of the instrument
\citep{feddersen1996}.

\paragraph{The status signal.}
The filer is optimizing and targets the highest-ranked opponent, which is the
premise Proposition~\ref{prop:foil} takes from \citet{nakaguma2025}. So the
identity of the target is informative about who is ahead, and its precision is
high, because the filer knows the race. A court case also draws local
attention, and revealing which candidates are first and second is known to
move votes \citep{anagol2016}.

Litigation is thus precise about position and imprecise about quality.

\section{Consolidation}
\label{sec:consolidation}

\begin{proposition}[Consolidation]
\label{prop:consolidation}
As $\lambda$ rises, rank moves toward common knowledge while valence beliefs
barely update. In a coordination equilibrium over the two leading candidates
\citep{myerson1993,fey1997,cox1997}, support flows from the runner-up to the
leader. The identity of the leader is unchanged and the size of the field is
unchanged.
\end{proposition}

The candidate-supply null is structural rather than incidental: the status
signal arrives after the registration deadline, so it cannot act on the
candidate set.

\begin{proposition}[Concentration indices are redundant]
\label{prop:concentration}
If the shift is a one-parameter reallocation between the top two, then every
summary concentration index -- the effective number of candidates, the
Herfindahl index, the combined top-two share -- is a monotone function of that
same parameter, and none of them carries a restriction independent of the
margin.
\end{proposition}

This is why concentration cannot identify which mechanism is at work. The
point is not that the concentration estimates are imprecise; it is that even a
precise one would be the margin restated. Pushing the estimated top-two
reallocation through each municipality's own 2024 share vector implies
$\Delta\mathrm{HHI} = +0.018$ and $\Delta\mathrm{ENP} = -0.053$, and the
separately estimated index regressions give $+0.023$ (0.013) and $-0.044$
(0.039). Both implied values sit inside the estimated intervals, so the index
regressions add no restriction to the margin result
(\texttt{docs/model/check\_proposition2.py}).

\section{Exit and the seat split}
\label{sec:seat}

\begin{proposition}[Exit]
\label{prop:exit}
A supporter of the runner-up who learns the race is not close faces a
collapsed pivot probability, so the choice becomes expressive: mark the leader
if $\hat\theta_1 + \varepsilon_{i1} \ge \bar b$, and otherwise mark blank.
Under compulsory voting the exit option is a blank or null vote rather than
abstention, so turnout is unmoved \citep{ghirardato2006}.
\end{proposition}

\begin{proposition}[The seat split]
\label{prop:seat}
The runner-up loses a similar mass whichever kind of seat is contested. What
differs is where that mass goes: to the leader where the revealed leader is a
fresh candidate, and off the valid ballot where the revealed leader is an
incumbent.
\end{proposition}

Two primitives differ across seat types, and both push the same way.

First, the prior precision $h_1$. In a contested seat the leader is the
incumbent, $h_1$ is high, $\hat\theta_1 \approx m_1$, and the type signal moves
nothing. In an open seat the leader is a fresh candidate with a diffuse
reputation and a genuinely unresolved rank, so the status signal has more to
resolve and coordination is stronger.

Second, selection on $\varepsilon_{i1}$. Supporting the challenger against a
sitting mayor is itself evidence of a settled negative view of that mayor, so
the distribution of $\varepsilon_{i1}$ among the runner-up's supporters is
first-order dominated in contested seats. A larger mass sits below $\bar b$
and exits. An open-seat leader has not been in office long enough to have
accumulated committed opponents.

\section{What the model does not do}
\label{sec:scope}

The model produces signs and orderings. It is not calibrated and not
structurally estimated. It does not separate a genuine allegation from a
persuasive weapon, and Section~\ref{sec:signals} is the reason that is a
property of the environment rather than a gap to be closed. It does not
endogenize the filer's decision: the attack direction is taken from
\citet{nakaguma2025} as a maintained assumption, because endogenizing it would
require exactly the selection the instrument exists to sidestep. It says
nothing about judges, and nothing about welfare. Runoff municipalities, where
a second round changes the coordination problem, are outside it.

Two conditions bound the mechanism. Coordination needs a small number of
focal contenders, so a genuinely three-way race weakens it. And the reasoning
in Section~\ref{sec:consolidation} applies to a single-winner plurality
contest; an open-list proportional council race offers no focal rank for the
status signal to reveal, which is the model's clearest out-of-sample
implication.

\bibliographystyle{apalike}
\bibliography{../../output/presentation/biblio}

\end{document}
